\documentclass[12pt]{ximera}
\typeout{Start loading xmPreamble.tex}%

% Add here extra macro's that are loaded automatically by all documents of claas 'ximera' or 'xourse' in this repo

%%
%%  Example:
%%
% \newcommand{\R}{\mathbb{R}}
\usepackage{tikz}
\usetikzlibrary{positioning, arrows.meta}
\usepackage{multicol}
\usepackage{enumitem}
\usepackage{caption}
\usepackage{amsmath}
\usepackage{amsthm}
\usepackage{fontenc}

%% ---------------------------------------------------------------
%% Shared worksheet macros.
%%
%% These came from the Summer 2026 worksheet set, where each worksheet carried its
%% own copy. They have to live here instead: a xourse inlines every activity into a
%% single document, so duplicate \newcommand definitions across activities collide,
%% and the xourse gobbles \input so a per-topic macro file never loads.
%%
%% They use \providecommand, NOT \newcommand. ximera.cls loads this file automatically,
%% and every activity also carries an explicit \typeout{Start loading xmPreamble.tex}%

% Add here extra macro's that are loaded automatically by all documents of claas 'ximera' or 'xourse' in this repo

%%
%%  Example:
%%
% \newcommand{\R}{\mathbb{R}}
\usepackage{tikz}
\usetikzlibrary{positioning, arrows.meta}
\usepackage{multicol}
\usepackage{enumitem}
\usepackage{caption}
\usepackage{amsmath}
\usepackage{amsthm}
\usepackage{fontenc}

%% ---------------------------------------------------------------
%% Shared worksheet macros.
%%
%% These came from the Summer 2026 worksheet set, where each worksheet carried its
%% own copy. They have to live here instead: a xourse inlines every activity into a
%% single document, so duplicate \newcommand definitions across activities collide,
%% and the xourse gobbles \input so a per-topic macro file never loads.
%%
%% They use \providecommand, NOT \newcommand. ximera.cls loads this file automatically,
%% and every activity also carries an explicit \typeout{Start loading xmPreamble.tex}%

% Add here extra macro's that are loaded automatically by all documents of claas 'ximera' or 'xourse' in this repo

%%
%%  Example:
%%
% \newcommand{\R}{\mathbb{R}}
\usepackage{tikz}
\usetikzlibrary{positioning, arrows.meta}
\usepackage{multicol}
\usepackage{enumitem}
\usepackage{caption}
\usepackage{amsmath}
\usepackage{amsthm}
\usepackage{fontenc}

%% ---------------------------------------------------------------
%% Shared worksheet macros.
%%
%% These came from the Summer 2026 worksheet set, where each worksheet carried its
%% own copy. They have to live here instead: a xourse inlines every activity into a
%% single document, so duplicate \newcommand definitions across activities collide,
%% and the xourse gobbles \input so a per-topic macro file never loads.
%%
%% They use \providecommand, NOT \newcommand. ximera.cls loads this file automatically,
%% and every activity also carries an explicit \input{../xmPreamble}, so this file is
%% read twice. That was harmless while it held only \usepackage lines (LaTeX skips a
%% package it has already loaded) but a second \newcommand is a hard error.
%%
%%   \blank[width]        fill-in-the-blank rule (default 2.5cm)
%%   \showwork            "show and explain your work" reminder
%%   \showworkline        ... plus which schedule line was used
%%   \showworkyear        ... plus which year's schedule was used
%%   \schedhead{...}      centred bold caption above a tiered-rate table
%%   \samesquares[scale]{left}{right}   two equal-area squares, labelled
%%   \samecubes[scale]{left}{right}     two equal-volume cubes, labelled
%% ---------------------------------------------------------------

\providecommand{\blank}[1][2.5cm]{\underline{\hspace{#1}}}
\providecommand{\showwork}{\textbf{REMEMBER TO SHOW AND EXPLAIN YOUR WORK.}}
\providecommand{\showworkline}{\textbf{REMEMBER TO SHOW AND EXPLAIN YOUR WORK.
  Explanation should include how and why you know which line of the
  schedule was used to calculate this amount.}}
\providecommand{\showworkyear}{\begin{sloppypar}\textbf{REMEMBER TO SHOW AND
  EXPLAIN YOUR WORK. Explanation should include how and why you know which
  line of the year's tax schedule was used to calculate this tax
  amount.}\end{sloppypar}}
\providecommand{\schedhead}[1]{\begin{center}\textbf{#1}\end{center}}

\providecommand{\samesquares}[3][1]{%
\begin{center}
{\footnotesize These squares are the \textbf{same size}.}

\vspace{1.5mm}
\begin{tikzpicture}[scale=#1,every node/.style={scale=#1}]
  \draw (0,0) rectangle (2.4,2.4);
  \node[anchor=north west] at (0.15,2.25) {Area:};
  \node[anchor=east]  at (-0.15,1.2) {#2};
  \node[anchor=north] at (1.2,-0.15) {#2};
  \begin{scope}[xshift=4.6cm]
    \draw (0,0) rectangle (2.4,2.4);
    \node[anchor=north west] at (0.15,2.25) {Area:};
    \node[anchor=east]  at (-0.15,1.2) {#3};
    \node[anchor=north] at (1.2,-0.15) {#3};
  \end{scope}
\end{tikzpicture}
\end{center}}

\providecommand{\samecubes}[3][1]{%
\begin{center}
{\footnotesize These cubes are the \textbf{same size}.}

\vspace{1.5mm}
\begin{tikzpicture}[scale=#1,every node/.style={scale=#1}]
  \foreach \dx in {0,5.2} {
    \begin{scope}[xshift=\dx cm]
      \draw (0,0) rectangle (2.0,2.0);
      \draw (0,2.0) -- (0.65,2.6) -- (2.65,2.6) -- (2.0,2.0);
      \fill[black!12] (2.0,0) -- (2.65,0.6) -- (2.65,2.6) -- (2.0,2.0) -- cycle;
      \draw (2.0,0) -- (2.65,0.6) -- (2.65,2.6) -- (2.0,2.0);
      \node[anchor=north west] at (0.1,1.9) {Volume:};
    \end{scope}
  }
  \node[anchor=east]  at (-0.15,1.0) {#2};
  \node[anchor=north] at (1.0,-0.15) {#2};
  \node[anchor=west]  at (2.25,0.4) {#2};
  \node[anchor=east]  at (5.05,1.0) {#3};
  \node[anchor=north] at (6.2,-0.15) {#3};
  \node[anchor=west]  at (7.45,0.4) {#3};
\end{tikzpicture}
\end{center}}

%% NOTE: do NOT add \usepackage to individual activity .tex files.
%% When a xourse inlines an activity it gobbles \documentclass and \input, but a bare
%% \usepackage still executes -- after \begin{document} -- and errors out.
%% All shared packages and macros belong here.
%%
%% ximera.cls already loads: amsmath, amssymb, amsthm, cancel, enumitem, graphicx,
%% hyperref, listings, pgfplots, tikz, url, xcolor. No need to re-load those.
, so this file is
%% read twice. That was harmless while it held only \usepackage lines (LaTeX skips a
%% package it has already loaded) but a second \newcommand is a hard error.
%%
%%   \blank[width]        fill-in-the-blank rule (default 2.5cm)
%%   \showwork            "show and explain your work" reminder
%%   \showworkline        ... plus which schedule line was used
%%   \showworkyear        ... plus which year's schedule was used
%%   \schedhead{...}      centred bold caption above a tiered-rate table
%%   \samesquares[scale]{left}{right}   two equal-area squares, labelled
%%   \samecubes[scale]{left}{right}     two equal-volume cubes, labelled
%% ---------------------------------------------------------------

\providecommand{\blank}[1][2.5cm]{\underline{\hspace{#1}}}
\providecommand{\showwork}{\textbf{REMEMBER TO SHOW AND EXPLAIN YOUR WORK.}}
\providecommand{\showworkline}{\textbf{REMEMBER TO SHOW AND EXPLAIN YOUR WORK.
  Explanation should include how and why you know which line of the
  schedule was used to calculate this amount.}}
\providecommand{\showworkyear}{\begin{sloppypar}\textbf{REMEMBER TO SHOW AND
  EXPLAIN YOUR WORK. Explanation should include how and why you know which
  line of the year's tax schedule was used to calculate this tax
  amount.}\end{sloppypar}}
\providecommand{\schedhead}[1]{\begin{center}\textbf{#1}\end{center}}

\providecommand{\samesquares}[3][1]{%
\begin{center}
{\footnotesize These squares are the \textbf{same size}.}

\vspace{1.5mm}
\begin{tikzpicture}[scale=#1,every node/.style={scale=#1}]
  \draw (0,0) rectangle (2.4,2.4);
  \node[anchor=north west] at (0.15,2.25) {Area:};
  \node[anchor=east]  at (-0.15,1.2) {#2};
  \node[anchor=north] at (1.2,-0.15) {#2};
  \begin{scope}[xshift=4.6cm]
    \draw (0,0) rectangle (2.4,2.4);
    \node[anchor=north west] at (0.15,2.25) {Area:};
    \node[anchor=east]  at (-0.15,1.2) {#3};
    \node[anchor=north] at (1.2,-0.15) {#3};
  \end{scope}
\end{tikzpicture}
\end{center}}

\providecommand{\samecubes}[3][1]{%
\begin{center}
{\footnotesize These cubes are the \textbf{same size}.}

\vspace{1.5mm}
\begin{tikzpicture}[scale=#1,every node/.style={scale=#1}]
  \foreach \dx in {0,5.2} {
    \begin{scope}[xshift=\dx cm]
      \draw (0,0) rectangle (2.0,2.0);
      \draw (0,2.0) -- (0.65,2.6) -- (2.65,2.6) -- (2.0,2.0);
      \fill[black!12] (2.0,0) -- (2.65,0.6) -- (2.65,2.6) -- (2.0,2.0) -- cycle;
      \draw (2.0,0) -- (2.65,0.6) -- (2.65,2.6) -- (2.0,2.0);
      \node[anchor=north west] at (0.1,1.9) {Volume:};
    \end{scope}
  }
  \node[anchor=east]  at (-0.15,1.0) {#2};
  \node[anchor=north] at (1.0,-0.15) {#2};
  \node[anchor=west]  at (2.25,0.4) {#2};
  \node[anchor=east]  at (5.05,1.0) {#3};
  \node[anchor=north] at (6.2,-0.15) {#3};
  \node[anchor=west]  at (7.45,0.4) {#3};
\end{tikzpicture}
\end{center}}

%% NOTE: do NOT add \usepackage to individual activity .tex files.
%% When a xourse inlines an activity it gobbles \documentclass and \input, but a bare
%% \usepackage still executes -- after \begin{document} -- and errors out.
%% All shared packages and macros belong here.
%%
%% ximera.cls already loads: amsmath, amssymb, amsthm, cancel, enumitem, graphicx,
%% hyperref, listings, pgfplots, tikz, url, xcolor. No need to re-load those.
, so this file is
%% read twice. That was harmless while it held only \usepackage lines (LaTeX skips a
%% package it has already loaded) but a second \newcommand is a hard error.
%%
%%   \blank[width]        fill-in-the-blank rule (default 2.5cm)
%%   \showwork            "show and explain your work" reminder
%%   \showworkline        ... plus which schedule line was used
%%   \showworkyear        ... plus which year's schedule was used
%%   \schedhead{...}      centred bold caption above a tiered-rate table
%%   \samesquares[scale]{left}{right}   two equal-area squares, labelled
%%   \samecubes[scale]{left}{right}     two equal-volume cubes, labelled
%% ---------------------------------------------------------------

\providecommand{\blank}[1][2.5cm]{\underline{\hspace{#1}}}
\providecommand{\showwork}{\textbf{REMEMBER TO SHOW AND EXPLAIN YOUR WORK.}}
\providecommand{\showworkline}{\textbf{REMEMBER TO SHOW AND EXPLAIN YOUR WORK.
  Explanation should include how and why you know which line of the
  schedule was used to calculate this amount.}}
\providecommand{\showworkyear}{\begin{sloppypar}\textbf{REMEMBER TO SHOW AND
  EXPLAIN YOUR WORK. Explanation should include how and why you know which
  line of the year's tax schedule was used to calculate this tax
  amount.}\end{sloppypar}}
\providecommand{\schedhead}[1]{\begin{center}\textbf{#1}\end{center}}

\providecommand{\samesquares}[3][1]{%
\begin{center}
{\footnotesize These squares are the \textbf{same size}.}

\vspace{1.5mm}
\begin{tikzpicture}[scale=#1,every node/.style={scale=#1}]
  \draw (0,0) rectangle (2.4,2.4);
  \node[anchor=north west] at (0.15,2.25) {Area:};
  \node[anchor=east]  at (-0.15,1.2) {#2};
  \node[anchor=north] at (1.2,-0.15) {#2};
  \begin{scope}[xshift=4.6cm]
    \draw (0,0) rectangle (2.4,2.4);
    \node[anchor=north west] at (0.15,2.25) {Area:};
    \node[anchor=east]  at (-0.15,1.2) {#3};
    \node[anchor=north] at (1.2,-0.15) {#3};
  \end{scope}
\end{tikzpicture}
\end{center}}

\providecommand{\samecubes}[3][1]{%
\begin{center}
{\footnotesize These cubes are the \textbf{same size}.}

\vspace{1.5mm}
\begin{tikzpicture}[scale=#1,every node/.style={scale=#1}]
  \foreach \dx in {0,5.2} {
    \begin{scope}[xshift=\dx cm]
      \draw (0,0) rectangle (2.0,2.0);
      \draw (0,2.0) -- (0.65,2.6) -- (2.65,2.6) -- (2.0,2.0);
      \fill[black!12] (2.0,0) -- (2.65,0.6) -- (2.65,2.6) -- (2.0,2.0) -- cycle;
      \draw (2.0,0) -- (2.65,0.6) -- (2.65,2.6) -- (2.0,2.0);
      \node[anchor=north west] at (0.1,1.9) {Volume:};
    \end{scope}
  }
  \node[anchor=east]  at (-0.15,1.0) {#2};
  \node[anchor=north] at (1.0,-0.15) {#2};
  \node[anchor=west]  at (2.25,0.4) {#2};
  \node[anchor=east]  at (5.05,1.0) {#3};
  \node[anchor=north] at (6.2,-0.15) {#3};
  \node[anchor=west]  at (7.45,0.4) {#3};
\end{tikzpicture}
\end{center}}

%% NOTE: do NOT add \usepackage to individual activity .tex files.
%% When a xourse inlines an activity it gobbles \documentclass and \input, but a bare
%% \usepackage still executes -- after \begin{document} -- and errors out.
%% All shared packages and macros belong here.
%%
%% ximera.cls already loads: amsmath, amssymb, amsthm, cancel, enumitem, graphicx,
%% hyperref, listings, pgfplots, tikz, url, xcolor. No need to re-load those.

\title{In-Class: Working with Data}
\begin{document}
\begin{abstract}
Mean, median, and mode; why a mean of category means is not the overall mean; what a
uniform percent raise does to each; reading percentiles; and using standard deviation
and the empirical rule on normally distributed data.
\end{abstract}
\maketitle

\section*{Mean, Median, and Mode}

\begin{example}
Consider the data set $2, 4, 6, 10, 13$.

\textbf{Mean (average).}
\[
\frac{\text{sum of values}}{\text{number of values}}
= \frac{2 + 4 + 6 + 10 + 13}{5} = \frac{35}{5} = 7.
\]
Note that rearranging this gives a fact we will use later:
$\text{mean} \cdot (\text{\# of values}) = \text{sum of the data points}$.

\textbf{Median.} First order the data from smallest to largest. If there is an odd
number of data points, the median is the middle value; if there is an even number,
there are two middle values, and these are averaged.

Here the data is already ordered and there are an odd number of points, so the median
is the middle value, $6$. Notice the median and the mean are \emph{not} equal.
\end{example}

\begin{example}
Now consider $2, 4, 6, 10, 13, 20$. The data is ordered and there are an even number
of points, so the two middle values are $6$ and $10$:
\[
\text{median} = \frac{6 + 10}{2} = \frac{16}{2} = 8.
\]
\end{example}

\subsection*{The Wing Aero data}

Let's return to the Wing Aero data from the pre-class work. Many companies keep
salary information private, so this is an imagined example of a fictional company.

% TODO: the Wing Aero salary table image is a broken-image placeholder in the source
% PDF (page 2) -- the underlying salary data is not in the document. The values below
% come from the instructor's worked answers. Request the original table; supplying it
% as real data rather than a picture would be better for accessibility anyway.

The company has $30$ employees, in four categories: $9$ workers, $9$ supervisors,
$8$ managers, and $4$ executives.

\begin{problem}
Record your information from your pre-class work.

Mean salary of all employees: \$$\answer[tolerance=0.05]{77.17}$k

Mean salary by category --- Worker: \$$\answer[tolerance=0.05]{25.44}$k \quad
Supervisor: \$$\answer[tolerance=0.05]{40.44}$k \quad
Manager: \$$\answer[tolerance=0.05]{79.13}$k \quad
Executive: \$$\answer[tolerance=0.05]{272.25}$k

\begin{solution}
Each mean is the sum of that group's salaries divided by how many are in the group:
\[
\frac{\text{sum of all salaries}}{30} = \$77.17\text{k},
\qquad
\frac{\text{sum of all worker salaries}}{9} = \$25.44\text{k}.
\]
\end{solution}
\end{problem}

\begin{problem}
How could we use the mean category salaries and the number of employees in each
category to determine the mean salary of all employees? Notice that you cannot simply
add up the four category means and divide by 4 --- why doesn't that work?

\begin{freeResponse}
\end{freeResponse}

\begin{solution}
Rearranging the definition of a mean, $9 \cdot \$25.44\text{k}$ is the sum of all the
worker salaries. Doing that for each category recovers the total payroll, and then
\[
\frac{9(\$25.44\text{k}) + 9(\$40.44\text{k}) + 8(\$79.13\text{k}) + 4(\$272.25\text{k})}{30}
= \$77.17\text{k},
\]
the mean for all employees.

Dividing the four category means by 4 doesn't work because it treats all four
categories as equally common. They are not --- there are 9 workers but only 4
executives. Averaging the means as if each carried the same weight would let the 4
executives' very high mean pull the result far too high. This is called a
\emph{weighted} mean: each category mean must be weighted by how many people it
describes.
\end{solution}
\end{problem}

\begin{problem}
Record your information from your pre-class work.

Median of all employees: \$$\answer[tolerance=0.05]{42.5}$k

Median by category --- Worker: \$$\answer[tolerance=0.05]{25}$k \quad
Supervisor: \$$\answer[tolerance=0.05]{42}$k \quad
Manager: \$$\answer[tolerance=0.05]{79.5}$k \quad
Executive: \$$\answer[tolerance=0.05]{270}$k
\end{problem}

\begin{problem}
\textbf{DO NOT RECALCULATE.} If everyone gets a $2.75\%$ raise, what happens to the
\emph{median} of each category? Will it go up or down, and by how much as a percent?

The median goes up by $\answer[tolerance=0.01]{2.75}$\%.

\begin{solution}
Suppose we have data $a, b, c, d, e$ ordered from smallest to largest. The median is
$c$. Increasing every value by $2.75\%$ gives
\[
1.0275a,\ 1.0275b,\ 1.0275c,\ 1.0275d,\ 1.0275e,
\]
which is still in order, so the new median is $1.0275c$ --- the median increased by
$2.75\%$.

The even case works the same way. For $a, b, c, d, e, f$ the median is
$\frac{c+d}{2}$, and after the raise
\[
\frac{1.0275c + 1.0275d}{2} = 1.0275 \cdot \frac{c+d}{2}.
\]
So if every data point increases by $2.75\%$, the median increases by $2.75\%$.
\end{solution}
\end{problem}

\begin{problem}
\textbf{DO NOT RECALCULATE.} If everyone gets a $2.75\%$ raise, what happens to the
\emph{mean} of each category?

The mean goes up by $\answer[tolerance=0.01]{2.75}$\%.

\begin{solution}
For data $a, b, c, d, e$ the mean is $\frac{a+b+c+d+e}{5}$. Increasing each value by
$2.75\%$ means replacing $a$ by $a + 0.0275a = 1.0275a$, and so on. The new mean is
\[
\frac{1.0275a + 1.0275b + 1.0275c + 1.0275d + 1.0275e}{5}
= \frac{1.0275(a+b+c+d+e)}{5}
= 1.0275 \cdot \text{original mean}.
\]
So the mean also increases by $2.75\%$. The common factor simply pulls out front ---
which is why neither statistic required any recalculation.
\end{solution}
\end{problem}

\section*{Percentiles}

Below is information about 1-year-old males' weights.

% TODO: the percentile table was a broken-image placeholder in the source PDF
% (pages 6 and 10). The rows below were reconstructed from the percentile/weight
% pairs quoted in the instructor's worked answers, so the table is PARTIAL -- it is
% every pair the worked solutions happen to mention, not the full original. Request
% the complete table.

\begin{center}
\begin{tabular}{|c|c|}
\hline
Weight & Percentile \\
\hline
19.5 lbs & 8.2nd \\
\hline
20.5 lbs & 17.5th \\
\hline
22.5 lbs & 46.5th \\
\hline
23.5 lbs & 61.8th \\
\hline
24.5 lbs & 74.9th \\
\hline
25.5 lbs & 84.8th \\
\hline
\end{tabular}
\end{center}

\begin{problem}
$25.5$ lbs corresponds to the 84.8th percentile. What does that mean, and what
percent of male 1-year-olds weigh \emph{more} than $25.5$ lbs?

$\answer[tolerance=0.05]{15.2}$\% weigh more than 25.5 lbs.

\begin{solution}
Being at the 84.8th percentile means $84.8\%$ of male 1-year-olds weigh the same or
less than $25.5$ lbs. Equivalently, since $100 - 84.8 = 15.2$, about $15.2\%$ weigh
more.
\end{solution}
\end{problem}

\begin{problem}
Remind yourself of the definition of the median. Based on that definition, what
percentile is the median for \emph{any} data set?

The $\answer{50}$th percentile.

\begin{solution}
The median is the middle value once the data is ordered, so half the data is the same
or less than the median. That is exactly the 50th percentile.
\end{solution}
\end{problem}

\begin{problem}
Using your answer above, estimate the median weight from the table. How did you
determine it?

About $\answer[tolerance=0.3]{22.7}$ lbs.

\begin{freeResponse}
\end{freeResponse}

\begin{solution}
The table does not list a value for the 50th percentile. It lists the 46.5th
percentile as $22.5$ lbs and the 61.8th as $23.5$ lbs, with nothing in between, so
the 50th percentile is somewhere between $22.5$ and $23.5$ lbs. Since $46.5$ is much
closer to $50$ than $61.8$ is, we estimate the 50th percentile at about $22.7$ lbs.
\end{solution}
\end{problem}

\section*{Standard Deviation}

Here is the image about standard deviation from the pre-class work.

\begin{center}
\includegraphics[width=4in]{workingWithData/normal-curve-percentiles}
\end{center}

\begin{problem}
If about $68\%$ of the data is within 1 standard deviation of the mean, what percent
lies \emph{outside} one standard deviation?
\begin{center}
$\answer{32}$\%
\end{center}

\begin{prompt}
Since the data within 1 standard deviation is the middle $68\%$, what percent of the
data is \emph{above} one standard deviation above the mean?
$\answer{16}$\%

And what percent is \emph{below} one standard deviation below the mean?
$\answer{16}$\%

\begin{solution}
$100\% - 68\% = 32\%$ lies outside. By symmetry that 32\% splits evenly between the
two tails, so $16\%$ in each.
\end{solution}
\end{prompt}

\begin{prompt}
So an interval showing one standard deviation either side of the mean starts around
the $\answer{16}$th percentile and ends around the $\answer{84}$th percentile.

\begin{solution}
The lower end has $16\%$ below it, putting it at the 16th percentile. The upper end
has $16\%$ above it, so $100 - 16 = 84\%$ at or below it --- the 84th percentile.
Equivalently $50\% - 34\% = 16\%$ and $50\% + 34\% = 84\%$.
\end{solution}
\end{prompt}
\end{problem}

\begin{problem}
Looking at the baby weight table above, the middle $68\%$ of 1-year-old males'
weights fall between about $\answer[tolerance=0.3]{20.3}$ lbs and
$\answer[tolerance=0.3]{25.5}$ lbs. Assume the weights are normally distributed.

\begin{hint}
You want the 16th and 84th percentiles. Neither is listed exactly --- estimate each
from the nearest listed percentiles, as you did for the median.
\end{hint}

\begin{solution}
The 16th percentile is not listed, but the 8.2nd percentile is $19.5$ lbs and the
17.5th is $20.5$ lbs, so the 16th must be between those weights. Since $17.5$ is much
closer to $16$ than $8.2$ is, estimate the 16th percentile at about $20.3$ lbs.

The 84th percentile is not listed either, but the 74.9th is $24.5$ lbs and the 84.8th
is $25.5$ lbs, so the 84th lies between them. Since $84.8$ is very close to $84$,
estimate the 84th percentile at about $25.5$ lbs.

So the middle $68\%$ of 1-year-old males' weights fall between about $20.3$ lbs and
$25.5$ lbs.
\end{solution}
\end{problem}

\subsection*{Further practice}

\begin{problem}
The score on a standardized test for a certain year is modeled using a normal
distribution. The mean is $76.5$ points and the standard deviation is $5.2$ points.
In the figure, $V$ is a number along the axis under the highest part of the curve,
and $U$ and $W$ are each the same distance away from $V$.

% TODO: BLOCKED -- the normal-distribution figure for this problem is a broken-image
% placeholder in the source PDF (page 11). Without it we cannot tell how much of the
% region is shaded, so part (a) is unanswerable and part (b) depends on it. Both are
% left ungraded pending the original figure.

\begin{prompt}
\textbf{(a)} Based on the empirical rule, how much of the data is in the shaded
region of the figure?
\begin{freeResponse}
\end{freeResponse}
\end{prompt}

\begin{prompt}
\textbf{(b)} Label $V$, $U$, and $W$ with their numerical values, based on the mean,
the standard deviation, and your answer to part (a).
\begin{freeResponse}
\end{freeResponse}

\begin{solution}
$V$ is under the highest part of the curve, so $V = 76.5$, the mean, whatever the
shading turns out to be.

For $U$ and $W$, work outward from the mean in steps of $5.2$: one standard deviation
gives $71.3$ and $81.7$ (containing $68\%$), two gives $66.1$ and $86.9$ ($95\%$),
three gives $60.9$ and $92.1$ ($99.7\%$). Which pair is correct depends on the shaded
region in part (a).
\end{solution}
\end{prompt}
\end{problem}

\begin{problem}
Suppose the mean systolic blood pressure for women over age seventy is $134$ mmHg,
with a standard deviation of $9$ mmHg, and that blood pressures are normally
distributed.

\begin{prompt}
\textbf{(a)} Approximately $95\%$ of women over seventy have blood pressures between
$\answer{116}$ mmHg and $\answer{152}$ mmHg.

\begin{solution}
The empirical rule puts $95\%$ of the data within \emph{two} standard deviations of
the mean. Two standard deviations is $2 \cdot 9 = 18$ mmHg, so
\[
134 - 18 = 116 \text{ mmHg}, \qquad 134 + 18 = 152 \text{ mmHg}.
\]
\end{solution}
\end{prompt}

\begin{prompt}
\textbf{(b)} Approximately $\answer[tolerance=0.2]{99.7}$\% of women over seventy
have blood pressures between $107$ mmHg and $161$ mmHg.

\begin{hint}
How many standard deviations from the mean is $107$? And $161$?
\end{hint}

\begin{solution}
Here we work backwards from the endpoints. Both are $27$ mmHg from the mean, and
\[
\frac{27}{9} = 3,
\]
so the interval is the mean plus or minus \emph{three} standard deviations. By the
empirical rule that captures about $99.7\%$ of the data.
\end{solution}
\end{prompt}
\end{problem}

\end{document}
